% Drawing 69/01-1 -- folder 69, batch 01, archive page 2 (PDF page 3).
% Transcription: en marge gauche, un triangle de sommets $w$ (en haut), $u$ et $v$ (en bas), et un point $x$ relié à $u$ par un trait
% Kind: triangle with a pendant edge (graph). Method 3 (encoded, rendered from the data). Status: ok.
% Notes: Triangle w (top), u, v (base), and x joined to u by a long stroke running to the left. The figure stands in the text column (left part), not strictly in the margin. Encoded as a graph (polyhedron/1 with one triangular face and an extra edge).
% Coordinates: 1 unit = 1 mm = 1/100 of the crop width (bbox in page fractions: [0.04, 0.6, 0.4, 0.75]).
% Redrawn by Claude (claude-opus-5-5), issue #73, 2026-11-07; not checked by a human; not compiled (no TeX in the session).
% Rights: reproduces part of the work -- behind RENDER_TRANSCRIPTS, never published by itself.
\documentclass[tikz,border=3pt]{standalone}
\usepackage{amsmath,amssymb}
\usetikzlibrary{patterns,arrows.meta}
\tikzset{pencil/.style={black!45}, struck/.style={black!35}, mark/.style={circle,draw,inner sep=1.3pt}, pt/.style={circle,fill,inner sep=0.9pt}, opt/.style={circle,draw,fill=white,inner sep=1.1pt}}
\begin{document}
\begin{tikzpicture}[x=1mm,y=1mm,line cap=round,line join=round,font=\scriptsize]
\draw (54,-22) -- (46,-30) -- (60,-31) -- cycle;
\draw (23,-37) -- (46,-30);
\node[pt] at (54,-22) {};
\node[pt] at (46,-30) {};
\node[pt] at (60,-31) {};
\node[pt] at (23,-37) {};
\node[inner sep=1pt] at (54,-17) {\scriptsize $w$};
\node[inner sep=1pt] at (43,-27) {\scriptsize $u$};
\node[inner sep=1pt] at (64,-33) {\scriptsize $v$};
\node[inner sep=1pt] at (20,-34) {\scriptsize $x$};
\end{tikzpicture}
\end{document}
